\section{Introduction}
\label{sec:introduction}

Cardiovascular pathologies, encompassing coronary artery disease, heart failure, cardiac arrhythmias, and acute myocardial infarction, represent the largest contributor to premature global morbidity and mortality \cite{Wagner2020}. Early identification and prompt electrophysiological intervention are essential for mitigating sudden cardiac death, particularly for occult arrhythmogenic syndromes such as paroxysmal atrial fibrillation and non-sustained ventricular tachycardia. 

Standard clinical diagnostics rely on the static 12-lead electrocardiogram (ECG) recorded under stationary conditions. While the 12-lead ECG offers high spatial resolution by capturing cardiac electrical vectors across frontal and horizontal planes, its episodic nature poses a fundamental limitation: transient electrophysiological abnormalities occurring outside the 10-second recording window routinely escape detection. To bridge this surveillance gap, ambulatory continuous monitors such as Holter monitors and adhesive chest patches have gained clinical prominence. Nevertheless, conventional Holter systems require multi-point chest lead placement, hydrogel electrodes that cause contact dermatitis over extended wear, and conspicuous wiring that impedes daily physical activity, resulting in high rates of patient non-compliance \cite{Goverdovsky2017}.

Recent advances in wearable biomedical engineering have highlighted the human ear as an exceptional anatomical site for continuous, unobtrusive physiological sensing \cite{vonRosenberg2017}. The external auditory canal and periauricular cartilage exhibit high biomechanical stability, reduced susceptibility to gross skeletal muscle movement during sedentary tasks, and dense vascularization via branches of the superficial temporal and posterior auricular arteries \cite{Allen2007, Celik2017}. Hearable devices that sit stably within the concha can acquire both cardiac biopotentials (auricular ECG) and microvascular volume changes (photoplethysmography, PPG) without the stigma or cutaneous discomfort associated with thoracic instrumentation.

However, translating raw auricular biopotentials into actionable, hospital-grade clinical inferences entails significant engineering and computational hurdles:
\begin{enumerate}
    \item \textbf{Low Signal-to-Noise Ratio (SNR) and Vector Attenuation:} Auricular cardiac potentials originate at a substantial physical distance from the cardiac dipole, resulting in signal amplitudes typically on the order of $20$--$100\,\mu\text{V}$, compared to $1$--$2\,\text{mV}$ observed in precordial chest leads. Consequently, auricular signals are highly vulnerable to baseline wander, electromyographic (EMG) noise from mastication or facial expression, and ambient electrical hum \cite{Clifford2012}.
    \item \textbf{The 12-Lead vs. Single-Lead Inductive Mismatch:} An architectural fallacy prevalent in recent wearable machine learning literature is the artificial replication of single-lead inputs across 12 channels to feed legacy convolutional neural networks trained on standard clinical datasets. In true human electrophysiology, leads aVR, aVL, aVF, and the precordial leads V1--V6 exhibit strict anatomical vector orientations. Cloning Lead-II signals destroys spatial vector coherence, inducing catastrophic diagnostic failures in clinical entities that rely on precordial voltage criteria, such as Ventricular Hypertrophy (HYP).
    \item \textbf{Synthetic Benchmark Traps and Class Imbalance:} Many published wearable models report inflated classification metrics by training on oversimplified, synthetic Gaussian isotropic clusters rather than real patient cohorts. When exposed to heterogeneous clinical populations characterized by severe class imbalance, such models fail in clinical triage.
    \item \textbf{Lack of Calibrated Uncertainty Quantification:} Conventional deep learning classifiers output overconfident softmax probabilities that do not reflect predictive epistemic uncertainty, creating hazards in autonomous emergency medical routing.
\end{enumerate}

To systematically resolve these challenges, this paper presents \textit{Ear to Heart}, a holistically integrated hardware, signal processing, and deep learning framework for continuous auricular cardiopulmonary assessment. The major contributions of this work are summarized as follows:
\begin{itemize}
    \item \textbf{Dual-Modal Auricular Hardware Front-End:} We design and validate an ear-worn sensor module integrating a dry-contact AD8232 ECG analog front-end (CMRR $> 80$\,dB) across a tragus-earlobe-canal montage and a MAX30102 dual-wavelength (660/880\,nm) PPG optical sensor, interfaced to an ultra-low-power ESP32 MCU streaming continuous 250\,Hz telemetry over BLE 5.0.
    \item \textbf{Real-Time Dual SQI and Adaptive Conditioning:} We implement a low-latency edge pipeline comprising zero-phase Butterworth bandpass filtering (0.5--40\,Hz), morphological baseline drift elimination, and multi-parameter composite Signal Quality Indices (SQI) for both ECG and PPG streams.
    \item \textbf{Native Single-Lead SE-ResNet-1D + BiGRU Architecture:} We formulate a deep neural network tailored specifically for single-lead auricular morphology, incorporating 1D Squeeze-and-Excitation channel attention and Bidirectional GRUs with temporal attention. Evaluated across the benchmark inter-patient partition of the PhysioNet MIT-BIH database (51,002 train beats and 47,432 test beats), the model achieves 98.02\% accuracy on Ventricular Ectopic Beats (PVCs).
    \item \textbf{Resolution of the Hypertrophy Collapse:} Across all 21,799 authentic PTB-XL single-lead records, our native single-lead formulation resolves the sensitivity collapse observed under cloned inputs, achieving 83.60\% Macro ROC-AUC.
    \item \textbf{Multi-Center Hemodynamic Stacking on 113,449 Patients:} We validate an end-to-end multimodal stacking ensemble across four major clinical cohorts (European CVD 70k, Heart Failure, Stroke Registry, Framingham), achieving 95.76\% ROC-AUC on heart failure stratification and 80.60\% on population CVD.
    \item \textbf{Interpretable Attribution and Conformal Safety:} We incorporate 1D Grad-CAM saliency mapping to visually confirm electrophysiological feature grounding alongside Inductive Split-Conformal Prediction providing provable 95\% marginal coverage guarantees.
\end{itemize}
